% Einheit 3 – Krümmung und Wendepunkte
\setcounter{aufgabe}{22}
\einheitenkopf[e3][3 Krümmung und Wendepunkte]{Einheit 3 von 5 · Krümmung und Wendepunkte}
\zweigzeile{Hier lernst du, mit der zweiten Ableitung Krümmung und Wendepunkte zu bestimmen und an einer gegebenen Stelle nachzuweisen · kennst du seit Q1 · Abitur GK · baut auf: Extrempunkte (Einheit 2), Gleichungen lösen, Ableitungen bilden (Nr. 2–5)}

\verfahren{Wendepunkte berechnen}
\begin{aufgabe}{Ich kann Wendepunkte berechnen.}
Bilde $f''$ und $f'''$.
\begin{teilezwei}
\tz $f(x) = x^3 - 9x^2$ & \tz $f(x) = -2x^3 + 5x^2$ \\
\tz $f(x) = -x^4 + 3x^2$ & \tz $f(x) = 0{,}25x^4 - x^3$ \\
\anweisung{Bestimme alle Wendepunkte.}
\tz $f(x) = x^3 - 6x^2 + 8x$ & \tz $f(x) = x^4 - 2x^3 + 1$ \\
\tz $f(x) = (x + 1)\cdot e^{-x}$ \quad (Abitur 2023 GK) & \\
\end{teilezwei}
\end{aufgabe}

\begin{aufgabe}{Ich kann angeben, wo ein Graph links- und wo er rechtsgekrümmt ist.}
Gib die Intervalle an, in denen der Graph von $f$ linksgekrümmt und in denen er rechtsgekrümmt ist.
\begin{teilezwei}
\tz $f(x) = x^3 + 3x^2 - 4$ & \tz $f(x) = -x^3 + 6x^2$ \\
\tz $f(x) = x^4 - 6x^2$ & \\
\end{teilezwei}
\end{aufgabe}

\verfahren{Wendepunkte nachweisen}
\begin{aufgabe}{Ich kann nachweisen, dass an einer gegebenen Stelle ein Wendepunkt liegt.}
Weise nach, dass der Graph von $f$ bei $x_0$ einen Wendepunkt hat.
\begin{teilezwei}
\tz $f(x) = x^3 + 6x^2 - 2$, \ $x_0 = -2$ & \tz $f(x) = -x^3 + 9x$, \ $x_0 = 0$ \\
\end{teilezwei}
\begin{teile}
\anweisung{Weise nach, dass $W$ ein Wendepunkt ist; bestätige auch den Funktionswert.}
\teil $f(x) = 0{,}5x^4 - 3x^2 + 2$, \ $W(1\,|\,{-}0{,}5)$
\anweisung{Bestimme den Wendepunkt und stelle die Gleichung der Wendetangente $w$ auf.}
\teil $f(x) = x^3 - 3x^2 + 4x$
\anweisung{Weise nach, dass $S$ ein Sattelpunkt ist.}
\teil $f(x) = x^4 + 2x^3 + 5$, \ $S(0\,|\,5)$
\anweisung{Begründe ohne Rechnung mithilfe einer Skizze.}
\teil Der Graph von $f$ mit $f(x) = (x - 3)\cdot e^{x}$ hat den Tiefpunkt $T(2\,|\,{-}e^2)$ und nähert sich für $x \to -\infty$ von unten der $x$-Achse. Begründe, dass der Graph von $f$ einen Wendepunkt hat. (Abitur 2021 GK)
\end{teile}
\end{aufgabe}

\begin{aufgabe}{Ich finde den Fehler bei Krümmung und Wendepunkten.}
Finde den Fehler, benenne ihn und korrigiere die Rechnung.
\begin{teile}
\teil Emma untersucht $f(x) = -x^3 + 3x^2$:
\rechnung{f''(x) &= -6x + 6 \\ f''(0) &= 6 > 0}
Emma schreibt: „Für $x < 1$ ist der Graph rechtsgekrümmt.“
\teil Paul untersucht $f(x) = x^3 + 3x^2$:
\rechnung{f''(x) &= 6x + 6 = 0 \quad\Longleftrightarrow\quad x = -1 \\ f'''(x) &= 6 \ne 0}
Paul schreibt: „Der Wendepunkt ist $W(-1\,|\,0)$.“
\end{teile}
\end{aufgabe}

\begin{aufgabe}{Ich kann begründen, wann ein Wendepunkt vorliegt.}
Entscheide ohne Rechnung und begründe.
\begin{teile}
\teil Für $f(x) = x^4 + 1$ gilt $f''(0) = 0$. Liegt bei $x = 0$ ein Wendepunkt?
\teil Warum durchsetzt die Wendetangente den Graphen im Wendepunkt, statt ihn dort nur zu berühren?
\end{teile}
\end{aufgabe}

\begin{aufgabe}{Ich kann Krümmung im Sachzusammenhang deuten.}
Eine Straße verläuft in der Draufsicht für $0 \le x \le 4$ entlang des Graphen von $f(x) = -0{,}1x^3 + 0{,}6x^2$ ($x$ und $f(x)$ in 100 m). Ein Auto fährt in Richtung wachsender $x$-Werte.
\begin{teile}
\teil Bestimme den Punkt, an dem die Straße von einer Linkskurve in eine Rechtskurve übergeht.
\teil Begründe, warum das Lenkrad in diesem Punkt einen Moment lang geradeaus steht.
\end{teile}
\end{aufgabe}
